\ssr{ЗАКЛЮЧЕНИЕ}

При выполнении лабораторной работы были выполнены следующие задачи:
\begin{enumerate}
	\item описан метод полного перебора, указаны его преимущество и недостатки, представлена схема и реализация алгоритма;
	\item описан муравьиный алгоритм без использования элитных муравьёв, указаны его преимущества и недостатки, представлена его схема и реализация;
	\item выполнена оценка трудоёмкости обоих алгоритмов на основе их схем;
	\item проведена параметризация муравьиного алгоритма;
	\item сформированы рекомендации по выбору значений параметров для заданного класса данных, состоящего как минимум из трёх полносвязных ориентированных графов;
	\item проведён сравнительный анализ рассмотренных методов с точки зрения качества решений, устойчивости и вычислительной сложности;
\end{enumerate}


Сравнительный анализ подтвердил, что метод полного перебора, несмотря на точность, неприменим на практике для крупных графов из-за экспоненциального роста сложности,
тогда как муравьиный алгоритм обеспечивает баланс между качеством решения и вычислительными затратами.
Таким образом, исследование подтвердило эффективность метаэвристических подходов для прикладных задач маршрутизации в асимметричных условиях.